\scanpage{56}

\[
\lambda_!\qquad A\qquad C
\]
\[
i_A(j_A(C))\longrightarrow C\qquad i^*(\Delta'_1)
\]
\[
i_Aj_a^o(C)\longrightarrow C^{\Delta^1}
\]
\figure{\begin{tikzpicture}\draw (0,0)--(.4,1.4);\end{tikzpicture}}

\diagram{\begin{tikzcd}
\mathcal M\arrow[r,shift left,"\lambda_!"]&\mathrm{Ch}_\bullet(Ab_k)\arrow[l,shift left,"\lambda^*"]
\end{tikzcd}}
\[
\lambda^*(\Theta_1)
\]
\diagram{\begin{tikzcd}[column sep=small]
&0\arrow[r,dash]&k\arrow[r,dash]&0\arrow[r,dash]&{}\\
\Theta_1&0\arrow[r,dash]&k\arrow[r]&k\arrow[u]\arrow[r,dash]&0\\
&&0\arrow[u]\arrow[r,dash]&k\arrow[u]\arrow[r,dash]&0
\end{tikzcd}}
\[
k\longleftarrow\Theta_1
\]
\figure{\begin{tikzpicture}
\draw (0,0)--(5,.3) .. controls (5.4,.3) and (5.2,-.2) .. (5.6,-.3) .. controls (7,-.8) and (9,-.7) .. (10,-.7);
\end{tikzpicture}}

\diagram{\begin{tikzcd}
\mathcal M\arrow[r,shift left,"\lambda_!"]&\mathrm{Ch}_\bullet(Ab_k)\arrow[l,shift left,"\lambda^*"]
\end{tikzcd}}
\[
\lambda^*(
\]
$\lambda^*$ transforme objets

\margin{
\diagram{\begin{tikzcd}M\arrow[r]&\mathrm{Ch}_\bullet(Ab_k)\end{tikzcd}}
\deleted{\unsure{$Ab_k$}}
\par\textcircled{2} (i)
$\lambda_!\lambda^*(l_\bullet)\longrightarrow l_\bullet$
\inserted{(\uline{asphérique} ?)}
\par\deleted{\unsure{$\lambda\otimes$}} (ii) $\rule{2cm}{.4pt}$
\par épi avec noyau quasi-nul
\par(iii) $\lambda_!,\lambda^*$ transforme \deleted{\unsure{équi.}} qu. iso en idem
et $\unsure{\mathrm{Hot}\mathbin{-}\mathcal M|\mathrm{Hot}_{\mathcal M}\longrightarrow\mathrm{Hot}}$ d'où transforme objet final en objet final
\par(iv) $\forall l_\bullet$ projectif,
$\deleted{\unsure{\lambda_!}}\lambda^*l_\bullet\mathbin{-}l_\bullet$ iso
\[
\left\{\begin{array}{l}iv_1\\iv_2\end{array}\right.
\]
\par(v) $\lambda_!(1_{\mathcal M})\longrightarrow k$ iso
(\unsure{Complexe} en $\rule{.6cm}{.4pt}$)
\par\deleted{\unsure{b)}} $\lambda^*$ transforme \inserted{\unsure{hom}}
\deleted{\unsure{nulles}} \unsure{acycliques} en idem
}

\rotatebox{90}{\begin{minipage}{1.35\linewidth}\ttfamily\small
\diagram{\begingroup\setlength{\arraycolsep}{\tabcolsep}$\begin{array}{rll}
\multicolumn{3}{l}{\text{\unsure{*******************************************************************************}}}\\
\text{*}&\text{JESNEWS DU 09/08/1982 C68 CADARACHE}&\text{*}\\
\multicolumn{3}{l}{\text{\unsure{*******************************************************************************}}}\\
\text{*}&\text{\quad DEPUIS LE 1/7/82 LES ALIAS DES COMMANDES LIBRARIAN}&\text{*}\\
\text{*}&\text{\qquad LIBGET (LG) ET LIBSAVE (LS) SONT SUPPRIMES,}&\text{*}\\
\text{*}&\text{\qquad POUR EVITER UNE POSSIBILITE DE CONFLIT AVEC}&\text{*}\\
\text{*}&\text{\qquad LE MEME ALIAS DES COMMANDES STANDARD IBM.}&\text{*}\\
\text{*}&\text{\quad ********}&\text{*}\\
\text{*}&\text{ATTENTION}&\text{*}\\
\text{*}&\text{\quad MISE EN SERVICE LE VENDREDI 30/7/82 A 8H00 DE LA DEUXIEME PHASE DE GAF}&\text{*}\\
\text{*}&\text{\qquad SUR 168 CADARACHE.}&\text{*}\\
\text{*}&\text{\qquad ELLE PERMET D'UTILISER DANS LA CDE. KEEP TSO ET DANS LA PROCEDURE KEEP}&\text{*}\\
\text{*}&\text{\qquad BATCH UNE POSSIBILITE D'AFFECTATION DE TYPE DE SAUVEGARDE AU NIVEAU DU}&\text{*}\\
\text{*}&\text{\qquad FICHIER.CETTE POSSIBILITE EST VIVEMENT CONSEILLEE.}&\text{*}\\
\text{*}&\text{\qquad POUR PLUS AMPLES INFORMATIONS VOIR NOTE TECHNIQUE CISI MVS140,LA NOTE}&\text{*}\\
\text{*}&\text{\qquad MVS 143 ETANT SUPPRIMEE.}&\text{*}\\
\text{*}&&\text{*}\\
\text{*}&\text{EXEMPLES D'UTILISATION:}&\text{*}\\
\text{*}&&\text{*}\\
\text{*}&\text{BATCH:\quad DSN=NOM.DU.FICHIER.A.SAUVEGARDER,DELAY=TYPE DE SAUVEGARDE}&\text{*}\\
\text{*}&&\text{*}\\
\text{*}&\text{TSO:}&\text{*}\\
\text{*}&\text{K DA('NOM DU FICHIER A SAUVEGARDER') DELAY(TYPE DE SAUVEGARDE)}&\text{*}\\
\text{*}&&\text{*}\\
\text{*}&\text{TYPES DE SAUVEGARDE:}&\text{*}\\
\text{*}&&\text{*}\\
\text{*}&\text{\quad W\qquad SAUVEGARDE HEBDO.}&\text{*}\\
\text{*}&\text{\quad D\qquad SAUVEGARDE JOURNALIERE.}&\text{*}\\
\text{*}&\text{\quad NOS\quad SAUVEGARDE MENSUELLE.}&\text{*}\\
\text{*}&\text{\quad W\qquad PARAMETRE PAR DEFAUT.}&\text{*}\\
\text{*}&&\text{*}\\
\multicolumn{2}{l}{\text{\unsure{*****************}}}&\text{*}\\
\text{*}&\text{\quad TSO\qquad *}&\text{*}\\
\multicolumn{2}{l}{\text{\unsure{*****************}}}&\text{*}\\
\text{*}&\text{\qquad LES NOUVELLES PROCEDURES PRT ET PUN,PLUS COMPLETES,SONT EN}&\text{*}\\
\text{*}&\text{\quad SERVICE A PARTIR DU 20/8/82.}&\text{*}\\
\text{*}&\text{\qquad VEUILLEZ CONSULTER:\qquad HELP PRT}&\text{*}\\
\text{*}&\text{\hspace{5cm}OU\qquad HELP PUN}&\text{*}\\
\multicolumn{3}{l}{\text{\unsure{*******************************************************************************}}}
\end{array}$\endgroup}

\diagram{\begingroup\setlength{\arraycolsep}{\tabcolsep}$\begin{array}{rll}
\multicolumn{3}{l}{\text{\unsure{******************************************************************}}}\\
\text{*}&&\text{*}\\
\text{*}&\text{\qquad - J E S N E W S\quad F R 1 -\qquad LE 19/08/82}&\text{*}\\
\text{*}&&\text{*}\\
\text{*}&\text{\unsure{----------------------------------------------------------------}}&\text{*}\\
\text{*}&&\text{*}\\
\text{*}&\text{DEPUIS LE LUNDI 12 JUILLET 82, LES CLASSES DE TRAVAUX D ET E}&\text{*}\\
\text{*}&\text{SERONT ETENDUES POUR UNE TAILLE MEMOIRE SUPERIEURE A 1024 K}&\text{*}\\
\text{*}&\text{SE REPORTER A LA NOTE TECHNIQUE MVS141.}&\text{*}\\
\text{*}&&\text{*}\\
\text{*}&\text{\unsure{----------------------------------------------------------------}}&\text{*}\\
\text{*}&\text{A PARTIR DU LUNDI 9 AOUT 1982 , IL SERA POSSIBLE D'IMPRIMER}&\text{*}\\
\text{*}&\text{SUR XEROX 9700 (DESTINATION SAC) DES FICHIERS UTILISANT LE}&\text{*}\\
\text{*}&\text{COUPLE DE JEUX DE CARACTERES : (GS15,GT15)}&\text{*}\\
\text{*}&&\text{*}\\
\multicolumn{3}{l}{\text{\unsure{******************************************************************}}}
\end{array}$\endgroup}
\end{minipage}}

\mspage{55}
